\documentclass[tikz,border=10pt]{standalone}
\usepackage{tikz}
\usepackage{tikz-3dplot}
\usepackage{amsmath}
\usetikzlibrary{arrows.meta}

\begin{document}

\tdplotsetmaincoords{65}{125}
\begin{tikzpicture}[tdplot_main_coords, scale=2.5]
    
    % Parâmetros do cubo
    \def\length{2}
    \pgfmathsetmacro{\half}{\length/2}
    
    % Desenhar arestas do cubo
    % Face frontal (z = half)
    \draw[thick] (-\half, -\half, \half) -- (\half, -\half, \half) -- (\half, \half, \half) -- (-\half, \half, \half) -- cycle;
    % Face traseira (z = -half)
    \draw[thick, dashed, gray!60] (-\half, -\half, -\half) -- (\half, -\half, -\half) -- (\half, \half, -\half) -- (-\half, \half, -\half) -- cycle;
    % Arestas conectando frente e trás
    \draw[thick] (-\half, -\half, -\half) -- (-\half, -\half, \half);
    \draw[thick] (\half, -\half, -\half) -- (\half, -\half, \half);
    \draw[thick] (\half, \half, -\half) -- (\half, \half, \half);
    \draw[thick, dashed, gray!60] (-\half, \half, -\half) -- (-\half, \half, \half);
    
    % Marcar os 8 vértices com coordenadas
    \foreach \x in {-1, 1} {
        \foreach \y in {-1, 1} {
            \foreach \z in {-1, 1} {
                \pgfmathsetmacro{\xcoord}{\x * \half}
                \pgfmathsetmacro{\ycoord}{\y * \half}
                \pgfmathsetmacro{\zcoord}{\z * \half}
                
                \fill (\xcoord, \ycoord, \zcoord) circle (0.03);
                
                % Posicionar labels de forma inteligente
                \ifnum\z=1
                    \ifnum\x=1
                        \ifnum\y=1
                            \node[above right, font=\large] at (\xcoord, \ycoord, \zcoord) 
                                {$\left(\frac{length}{2}, \frac{length}{2}, \frac{length}{2}\right)$};
                        \else
                            \node[below right, font=\large] at (\xcoord, \ycoord, \zcoord) 
                                {$\left(\frac{length}{2}, -\frac{length}{2}, \frac{length}{2}\right)$};
                        \fi
                    \else
                        \ifnum\y=1
                            \node[above left, font=\large] at (\xcoord, \ycoord, \zcoord) 
                                {$\left(-\frac{length}{2}, \frac{length}{2}, \frac{length}{2}\right)$};
                        \else
                            \node[below left, font=\large] at (\xcoord, \ycoord, \zcoord) 
                                {$\left(-\frac{length}{2}, -\frac{length}{2}, \frac{length}{2}\right)$};
                        \fi
                    \fi
                \fi
            }
        }
    }
    
    % Destacar uma face (face superior y = half)
    \draw[blue!60, line width=1.2pt, opacity=0.7] 
        (-\half, \half, -\half) -- (\half, \half, -\half) -- (\half, \half, \half) -- (-\half, \half, \half) -- cycle;
    
    % Eixos coordenados
    \draw[-Stealth, thick, blue!70!black] (0,0,0) -- (\half+0.7, 0, 0) node[right] {$z$};
    \draw[-Stealth, thick, red!70!black] (0,0,0) -- (0, \half+0.7, 0) node[above] {$x$};
    \draw[-Stealth, thick, green!70!black] (0,0,0) -- (0, 0, \half+0.7) node[above] {$y$};
    
    \fill (0,0,0) circle (0.03) node[below left, font=\large] {$(0,0,0)$};
    
    % Anotação
    \node[align=left, font=\large, draw, fill=white, rounded corners, opacity=0.9] at (-1.8, 0.5, 1.5) {
        6 faces: \\
        $XZ$: $y = \pm\frac{length}{2}$ \\
        $XY$: $z = \pm\frac{length}{2}$ \\
        $YZ$: $x = \pm\frac{length}{2}$
    };
    
\end{tikzpicture}

\end{document}